%%==========================================================================
\section{Conclusions and Future Work}
\label{sec:conclusion}
%%==========================================================================

This paper revisited the compact Tail Assignment Problem formulation of
\citet{khaled2018compact} and made three focused contributions. First, we
showed that the basic continuity and turn-time inequalities admit a
simultaneous-departure corner case and closed it with explicit non-overlap
constraints. Second, we proved that the paper's cumulative-flight-time
row can relax a maintenance interval incorrectly and that the
natural endpoint split is also insufficient. We provided an exact
cumulative-state correction and verified the failure against the production
constraint objects. Third, we developed a day-free event-based reformulation
for type-A maintenance and an ordered prefix-state implementation.

The event formulation attaches maintenance to its triggering flight and
propagates the flight-hour counter along the aircraft's selected route. In the
four-method horizon experiment on the corrected feasible family, all four
models were optimal. The three legacy variants had identical objectives,
and the event-based model used the same maintenance-cost accounting.
The explicit-state corrected legacy model was faster than the endpoint split
on every tested horizon. These results demonstrate the need to separate
feasibility correction from objective accounting; they do not support a
solver-independent speedup claim.

These computational results are deliberately scoped. They use deterministic
synthetic feasible instances, one instance per parameter cell, and a single
unrestricted CPLEX configuration. The results expose model-size and timing
differences among the four formulations on the tested A-only family;
they do not establish a universal speedup across solvers, instance
distributions, or larger benchmark sets.

\paragraph{Reproducibility.}
The project contains the model implementations, deterministic instance
generator, benchmark scripts, result tables, and plotting scripts needed to
reproduce the reported snapshots. The legacy, event, and ordered A-only
builders remain separately selectable so future comparisons can distinguish
model semantics from implementation choices.

\subsection*{Future directions}

\begin{itemize}
  \item \textbf{Larger A-only benchmark families:} repeat the horizon and
    fleet-size experiments over multiple seeds and heterogeneous timetable
    structures, reporting solver variability and confidence intervals after
    resolving the corrected legacy model's initial-state feasibility.
  \item \textbf{Integrated fleet assignment and routing:} extend the corrected
    compact model while retaining physically executable route constraints.
  \item \textbf{Robust and recoverable variants:} study uncertainty in flight
    times, disruptions, and maintenance availability under the exact state
    formulation.
  \item \textbf{Heuristic methods:} develop constructive and local-search
    methods for the corrected A-only model and compare them with the exact
    formulations on larger instances.
\end{itemize}
